%% Question bank chapter file - generated by bank_tools.py build (do not edit while the bank is still being built:
%% it is regenerated from the per-test data in Question Bank\data). Every question is one bankq block:
%%   \begin{bankq}[id=..., source={...}, paper=.., qno=.., type=..., subtopic={...}, answer={...}, answer_src=key|solved,
%%                 working={...}, difficulty=1-10, diff_parts={K S I C F}, diff_note={...},
%%                 group=<passage id or empty>, also_in={duplicate copies in other tests}]
%% difficulty = min(10, K+S+I+C+F-1): K concepts 1-3, S steps 1-4, I insight 0-2, C computation 0-2, F format 0-1
%%   ...question LaTeX, options, figures (\bfig / \bfigi -> img/)...
%%   \end{bankq}
%% Passages shared by several questions are bankpassage blocks; their questions carry group=<passage id>.
\documentclass[10pt]{article}
\usepackage[a4paper,margin=1.8cm]{geometry}
\usepackage[T1]{fontenc}
\usepackage{amsmath,amssymb,mathtools}
\usepackage{graphicx}
\usepackage[export]{adjustbox}
\usepackage{xcolor,needspace,keyval}
\usepackage{enumitem}
\usepackage{booktabs,array,tabularx}
\usepackage[hidelinks]{hyperref}
\setlength{\parindent}{0pt}
\setlength{\parskip}{3pt}
\makeatletter
\define@key{bq}{id}{\def\bq@id{#1}}\define@key{bq}{source}{\def\bq@src{#1}}\define@key{bq}{paper}{\def\bq@paper{#1}}
\define@key{bq}{qno}{\def\bq@qno{#1}}\define@key{bq}{type}{\def\bq@type{#1}}\define@key{bq}{subtopic}{\def\bq@sub{#1}}
\define@key{bq}{answer}{\def\bq@ans{#1}}\define@key{bq}{answer_src}{\def\bq@asrc{#1}}\define@key{bq}{working}{\def\bq@work{#1}}
\define@key{bq}{group}{\def\bq@group{#1}}\define@key{bq}{also_in}{\def\bq@also{#1}}
\define@key{bq}{difficulty}{\def\bq@diff{#1}}\define@key{bq}{diff_parts}{\def\bq@dparts{#1}}\define@key{bq}{diff_note}{\def\bq@dnote{#1}}
\newenvironment{bankq}[1][]{\def\bq@work{}\def\bq@also{}\def\bq@group{}\def\bq@diff{?}\def\bq@dparts{}\def\bq@dnote{}\setkeys{bq}{#1}%
  \par\needspace{4\baselineskip}\medskip\noindent{\small\textbf{\bq@id}\hfill\color{gray}[\bq@src, Paper \bq@paper, Q\bq@qno]
  \bq@type\ $\cdot$ \bq@sub\ $\cdot$ \textbf{D\bq@diff}}\par\nopagebreak\relax}
  {\par{\small\color{teal}Answer: \textbf{\bq@ans}\ifx\bq@work\@empty\else\ (solved: \bq@work)\fi
  \ifx\bq@also\@empty\else\par\color{gray}Also in: \bq@also\fi}\par}
\define@key{bp}{id}{\def\bp@id{#1}}\define@key{bp}{source}{\def\bp@src{#1}}\define@key{bp}{questions}{\def\bp@qs{#1}}
\newenvironment{bankpassage}[1][]{\setkeys{bp}{#1}\par\needspace{4\baselineskip}\bigskip\noindent
  {\small\textbf{Passage \bp@id}\hfill\color{gray}[\bp@src; questions \bp@qs]}\par\nopagebreak\relax\itshape}{\par}
\makeatother
\newcommand{\bfig}[2]{\par\smallskip{\centering\includegraphics[width=#1pt,max width=0.62\linewidth,max height=0.32\textheight]{img/#2}\par}\smallskip}
\newcommand{\bfigi}[2]{\includegraphics[width=#1pt,max width=\linewidth,max height=0.32\textheight,valign=c]{img/#2}}
\newcommand{\bmissing}[1]{\textit{[figure missing in the source conversion: \detokenize{#1}]}}
\newenvironment{qoptions}{\begin{itemize}[leftmargin=1.8em,itemsep=1pt,topsep=2pt,label={}]}{\end{itemize}}
\newcommand{\opt}[2]{\item[(#1)] #2}
\providecommand{\fwstart}{}\providecommand{\fwend}{}\providecommand{\qfullsplit}{}
\newcolumntype{L}[1]{>{\raggedright\arraybackslash}p{#1}}
\newcolumntype{C}[1]{>{\centering\arraybackslash}p{#1}}
\newcolumntype{R}[1]{>{\raggedleft\arraybackslash}p{#1}}
\providecommand{\sectionhdr}[1]{}\providecommand{\subsectionhdr}[1]{}


\begin{document}
{\Large\bfseries Environmental Chemistry --- Question Bank}\par\smallskip
{\small Chemistry chapter C34. Questions from all converted test papers, grouped by source series. Answers: printed key where it matches the paper, otherwise solved (answer\_src=solved, with working). Exact repeats across papers appear once, with also\_in listing the other copies.}\par\medskip\hrule

\section*{Score-II papers}
\begin{bankq}[id=SII07-P1-CHEM-S1-Q13, source={Score-II - Test 07 (10 Jul 2022)}, paper=1, qno=13, type=single, subtopic={Pollutants}, answer={D}, answer_src=key, difficulty=4, diff_parts={K2 S1 I1 C1 F0}, diff_note={concepts: standard problem in this chapter, pollutants | steps: apply the relevant results; compute}]
A sample of ground water has following contaminations :
\begin{center}
\begin{tabular}{|l|l|}\hline
Species & Concentration\\\hline
NO$_3^-$ & 35 ppm\\\hline
Pb & 20 ppb\\\hline
SO$_4^{2-}$ & 400 ppm\\\hline
Cd & 0.2 ppm\\\hline
\end{tabular}
\end{center}
This sample is unsuitable for drinking due to presence of :
\begin{qoptions}
\opt{A}{NO$_3^-$ ions}
\opt{B}{Pb}
\opt{C}{SO$_4^{2-}$ ions}
\opt{D}{Cd}
\end{qoptions}
\end{bankq}
\begin{bankq}[id=SII30-P1-CHEM-S1-Q18, source={Score-II - Test 30 (27 Mar 2024)}, paper=1, qno=18, type=single, subtopic={Pollutants}, answer={D}, answer_src=key, difficulty=4, diff_parts={K2 S1 I1 C1 F0}, diff_note={concepts: standard problem in this chapter, pollutants | steps: apply the relevant results; compute}]
Given below are two statements:\\
\textbf{Statement-I:} High concentration of strong nucleophilic reagent with secondary alkyl halides which do not have bulky substituents will follow S$_\mathrm{N}$2 mechanism.\\
\textbf{Statement-II:} A secondary alkyl halide when treated with a large excess of ethanol follows S$_\mathrm{N}$1 mechanism.\\
In the the light of the above statements, choose the most appropriate from the questions given below:
\begin{qoptions}
\opt{A}{Statement I is true but Statement II is false}
\opt{B}{Statement I is false but Statement II is true}
\opt{C}{Both statement I and Statement II are false}
\opt{D}{Both statement I and Statement II are true}
\end{qoptions}
\end{bankq}

\section*{Score Advanced papers}
\begin{bankq}[id=SA04-P2-CHEM-S1-Q02, source={Score Advanced - Test 04 (04 Jul 2021)}, paper=2, qno=2, type=single, subtopic={Pollutants}, answer={D}, answer_src=solved, working={CO2 does not produce photochemical smog}, difficulty=3, diff_parts={K1 S1 I1 C1 F0}, diff_note={concepts: smog precursors | steps: nitrogen oxides, hydrocarbons, acrolein and PAN take part; CO2 does not}]
Which of the following air pollutant is likely not to produce photo chemical smog?
\begin{qoptions}
\opt{A}{NO}
\opt{B}{CH$_2$CHCHO}
\opt{C}{\bfigi{55.4}{SA04_p2_chem_s1_q02_1.png}}
\opt{D}{CO$_2$}
\end{qoptions}
\end{bankq}
\begin{bankq}[id=SA04-P2-CHEM-S1-Q03, source={Score Advanced - Test 04 (04 Jul 2021)}, paper=2, qno=3, type=single, subtopic={Pollutants}, answer={B}, answer_src=solved, working={Statement B is false}, difficulty=4, diff_parts={K1 S2 I1 C1 F0}, diff_note={concepts: acid rain | steps: oxides of sulphur and nitrogen cause it; CO2 only lowers pH to 5.6; CO is not acidic}]
Which of the following is false regarding acid rain.
\begin{qoptions}
\opt{A}{Acid rain has pH below 5.6}
\opt{B}{CO$_2$ and CO are major contributors of acid rain}
\opt{C}{Acid rain washes away plant nutrients needed for their growth}
\opt{D}{Acid rain is responsible for the change of colour of Tajmahal.}
\end{qoptions}
\end{bankq}
\begin{bankq}[id=SA04-P2-CHEM-S1-Q17, source={Score Advanced - Test 04 (04 Jul 2021)}, paper=2, qno=17, type=single, subtopic={Ozone}, answer={C}, answer_src=solved, working={Decomposition of ozone releases energy, so the statement is incorrect}, difficulty=4, diff_parts={K2 S1 I1 C1 F0}, diff_note={concepts: incorrect statement, ozone | steps: ozone decomposition is exothermic}]
Which of the following statement is INCORRECT.
\begin{qoptions}
\opt{A}{In Haber's process iron oxide, potassium oxide and aluminium oxide are used as catalyst.}
\opt{B}{ClO$_2$ is used as bleaching agent}
\opt{C}{Decomposition of ozone is an endothermic process}
\opt{D}{Hydrolysis of XeF$_6$ is a non redox process}
\end{qoptions}
\end{bankq}
\begin{bankq}[id=SA05-P1-CHEM-S1-Q18, source={Score Advanced - Test 05 (15 Jul 2021)}, paper=1, qno=18, type=single, subtopic={Pollutants}, answer={C}, answer_src=key, difficulty=2, diff_parts={K1 S1 I0 C1 F0}, diff_note={concepts: photochemical smog components | steps: nitrogen oxide, PAN and acrolein}]
Air pollutants that produce photochemical smog are :
\begin{qoptions}
\opt{A}{Ozone, chloride and SO$_2$}
\opt{B}{O$_2$, Cl$_2$, HNO$_3$}
\opt{C}{Nitrogen oxide, PAN and Acrolein}
\opt{D}{CO$_2$, CO and SO$_2$}
\end{qoptions}
\end{bankq}

\section*{Nurture SRG papers}
\begin{bankq}[id=NSRG03-P1-CHEM-S2-Q09, source={Nurture SRG - Test 03 (04 Feb 2024)}, paper=1, qno=9, type=multi, subtopic={Pollutants}, answer={B,C,D}, answer_src=solved, working={Fluorapatite is harder (true); the lead limit is 50 ppb (not ppm), blue baby syndrome comes from nitrate (not sulphate) and nitrate does not cause constipation; the printed key lists A,B but A is a true statement}, difficulty=4, diff_parts={K1 S2 I1 C0 F1}, diff_note={concepts: statements on pollutants and water quality | steps: bone and tooth chemistry; limits in water; effects of sulphate and nitrate}]
Which one of the following is/are \textbf{INCORRECT} ?
\begin{qoptions}
\opt{A}{Fluorapatite is much harder than hydroxyapatite}
\opt{B}{Max. allowed limit of Pb in water is 50 ppm}
\opt{C}{Excess of sulphate in water causes blue baby syndrome}
\opt{D}{Excess of nitrate in water causes constipation}
\end{qoptions}
\end{bankq}
\begin{bankq}[id=NSRG03-P2-CHEM-S2-Q08, source={Nurture SRG - Test 03 (04 Feb 2024)}, paper=2, qno=8, type=multi, subtopic={Pollutants}, answer={A,C,D}, answer_src=solved, working={NO forms in lightning and is oxidised to the dioxide; ozone does not oxidise nitrogen to NO}, difficulty=4, diff_parts={K1 S2 I1 C0 F1}, diff_note={concepts: nitrogen oxides in the atmosphere | steps: formation of nitric oxide; oxidation to the dioxide; nitrous oxide}]
Select the correct options given below;
\begin{qoptions}
\opt{A}{NO is formed from $\mathrm{N_2}$ and $\mathrm{O_2}$, when lightening occurs at high altitude.}
\opt{B}{NO is formed when $\mathrm{N_2}$ is oxidised by $\mathrm{O_3}$}
\opt{C}{NO is oxidised to $\mathrm{NO_2}$ by $\mathrm{O_2}$ or by $\mathrm{O_3}$ in the stratosphere.}
\opt{D}{$\mathrm{N_2O}$ is formed by the reaction $\mathrm{N_2}$ with $\mathrm{O_3}$ or with atomic oxygen in the upper atmosphere and also by microbiological process.}
\end{qoptions}
\end{bankq}
\end{document}
